\documentclass[12pt]{article}

\usepackage[margin=1.10in]{geometry}
\usepackage{amsmath,amssymb,amsthm,mathtools}
\usepackage{array}
\usepackage{booktabs}
\usepackage{longtable}
\usepackage{hyperref}

\newtheorem{theorem}{Theorem}
\theoremstyle{definition}
\newtheorem{definition}{Definition}
\theoremstyle{remark}
\newtheorem{remark}{Remark}

\newcommand{\A}{\mathcal A}
\newcommand{\B}{\mathcal B}
\newcommand{\C}{\mathbb C}
\newcommand{\D}{\mathcal D}
\newcommand{\F}{\mathcal F}
\newcommand{\M}{\mathcal M}
\newcommand{\N}{\mathcal N}
\newcommand{\R}{\mathbb R}
\newcommand{\Z}{\mathcal Z}
\setlength{\emergencystretch}{3em}
\hypersetup{
  pdftitle={Support, Readout, and the Binary-Rank Obstruction},
  pdfsubject={Operator algebras, AQFT, operational reconstruction, and finite centers},
  pdfkeywords={AQFT, von Neumann algebras, finite centers, support projections, DHR sectors, locally covariant QFT, operator-algebra quantum error correction, black-hole entropy}
}

\title{Support, Readout, and the Binary-Rank Obstruction}
\author{}
\date{Submission-freeze package of June 12, 2026}

\begin{document}
\maketitle

\begin{abstract}
Operational reconstructions in algebraic quantum field theory and quantum
information reconstruct the algebraic object that is supplied: a represented
von Neumann algebra, a locally covariant functor, a sector category, a
correctable logical algebra, or a Type II algebra with a specified trace and
entropy normalization.  This note isolates a finite central obstruction to
turning such a reconstruction into a derivation of a binary microscopic
alphabet.  If the displayed data factor through a normal quotient
\(\pi_q:\widehat\A_q\to\A\), where
\(\widehat\A_q=\A\oplus(\N\bar\otimes\ell^\infty(\{1,\ldots,q\}))\), then the
same data are compatible with \(q=2\) and with \(q>2\) finite central
completions.  They therefore do not determine full ambient support,
faithfulness of the represented physical algebra, the physical/null/gauge/
inadmissible status of \(\ker\pi_q\), the finite central rank \(q\), or the
predicate \(q=2\).

The theorem also records the exact positive form of the missing hypotheses.
First, every non-null finite central direction must be accounted for by
available probes/readouts or by a null, gauge, or inadmissibility quotient:
for \(V_q=\ell^\infty_q/\C 1\), this is the rank equation
\(\operatorname{rank}R_q=\dim(V_q/N_q)\).  Second, even after support/readout
accountability, binary rank does not follow from entropy alone; every
\(q\ge2\) satisfies \(S_q=\log q=\Delta A_q/(4G\hbar)\) after choosing
\(\Delta A_q=4G\hbar\log q\).  A binary conclusion requires a further
finite-alphabet law, such as the least nontrivial closed quotient rule, or
direct \(q=2\) spectroscopy whose measurement channel is itself derived.  A
small Lean 4 file verifies the finite same-data/no-reconstruction kernel.
\end{abstract}

\paragraph{Keywords.}
Algebraic quantum field theory; von Neumann algebras; finite centers; support
projections; locally covariant quantum field theory; DHR sectors;
operator-algebra quantum error correction; black-hole spectroscopy.

\paragraph{Status.}
This is a submission-freeze package extracted from existing notes.  It does not
modify the frozen finite-center paper.  Its claim is intentionally negative and
conditional: quotient-visible data do not force a binary primitive.

\section{Question}

The intended target is not a new reconstruction theorem for AQFT, modular
theory, entanglement-wedge reconstruction, or gravitational crossed products.
Those positive theorems are accepted as inputs.  The question is a narrower
converse.  Suppose an operational experiment supplies normal state and effect
values, channels, modular data, locally covariant time-slice and relative
Cauchy evolution data, DHR/DR sector data, JLMS/OAQEC recovery data, continuum
limits, and entropy assignments, but all of the displayed data factor through a
represented quotient \(\A\).  Can those data alone prove that there is no
additional finite central support in a larger admitted algebra
\(\widehat\A_q\)?  Can they decide whether such a finite central kernel is
physical, null, gauge, or inadmissible?  Can they select the special finite
rank \(q=2\)?

The answer in this note is negative under explicit assumptions.  The same
answer is also the useful positive guide: a proof of binary rank must stop
using quotient-visible data and must exhibit the extra non-quotient input.  In
finite language, that input has two independent pieces.  The first is
support/readout accountability: all finite central directions must either be
separated by physical probes/readouts or be removed by a null/gauge/
inadmissibility rule.  The second is a primitive finite-alphabet rule: after a
finite center is accountable, entropy and ordinary yes/no tests still do not
tell us which quotient is counted as the elementary finite alphabet.

\section{Novelty Boundary}

The table gives the novelty check.  The note should be demoted if a
published theorem already states the joint criterion in the right form.

\begin{longtable}{p{0.28\linewidth}p{0.34\linewidth}p{0.30\linewidth}}
\toprule
Claim in this note & Nearest standard result & Honest boundary \\
\midrule
Operational data reconstruct only targets constant on data fibers. &
This is elementary descent for functions, and in operator algebras it appears
whenever a normal quotient or represented null space is fixed. &
No novelty is claimed for the fiber lemma.  The contribution is the finite
central application to support, readout, and binary rank. \\
\addlinespace
Quotient-visible locally covariant and relative-Cauchy data for a supplied
functor do not decide whether an invisible finite center has been admitted. &
Brunetti--Fredenhagen--Verch define locally covariant QFT and relative Cauchy
evolution; Fewster--Verch define dynamical locality and SPASs
\cite{BFV,FewsterVerch}. &
No new locally covariant theorem is claimed.  Full-theory dynamical locality,
source faithfulness, or indecomposability may be exactly the extra positive
assumption; quotient-visible data alone are weaker. \\
\addlinespace
DHR/DR sector data do not recreate central-character fibers omitted from the
supplied category. &
Doplicher--Haag--Roberts and Doplicher--Roberts reconstruct field algebras and
gauge groups from a supplied sector structure \cite{DHR1,DHR2,DR}. &
No new DR reconstruction theorem is claimed.  The statement is a warning about
which category was supplied. \\
\addlinespace
JLMS/OAQEC recovers a supplied logical algebra, not a larger finite central
completion. &
JLMS, Harlow, Dong--Harlow--Wall, and infinite-dimensional OAQEC work identify
relative entropy/recovery conditions for specified algebras
\cite{JLMS,HarlowRT,DHW,KangKolchmeyer}. &
The note does not dispute recovery.  It says maximality or full logical
algebra data is an additional input. \\
\addlinespace
Crossed-product/generalized-entropy data do not select \(q=2\) without finite
rank calibration. &
Witten, Chandrasekaran--Longo--Penington--Witten, and
Chandrasekaran--Penington--Witten construct Type II algebras and entropy
notions in gravitational settings \cite{WittenCrossed,CLPW,CPW}. &
The note does not rederive generalized entropy.  It isolates finite residuals
when a central register is outside the supplied algebra or trace calibration. \\
\addlinespace
Black-hole entropy does not by itself force binary rank. &
Bekenstein--Mukhanov black-hole spectroscopy relates area spacing, degeneracy,
and line emission \cite{BekensteinMukhanov}. &
A resolved \(\log 2\) line would validate binary empirically.  The note says
entropy alone is not that measurement. \\
\bottomrule
\end{longtable}

\section{Data and Targets}

Let \(\A\) be the represented algebra reconstructed from the displayed
operational experiment.  Let \(\N\) be a nonzero von Neumann algebra and, for
each \(q\ge2\), set
\[
  \widehat\A_q
  =
  \A\oplus\bigl(\N\bar\otimes \ell^\infty(\{1,\ldots,q\})\bigr).
\]
Let \(\pi_q:\widehat\A_q\to\A\) be the normal quotient onto the first summand.
The finite central projection
\[
  p_q=(0,1_{\N\bar\otimes\ell^\infty_q})
\]
is nonzero in \(\widehat\A_q\) and is killed by \(\pi_q\).

The displayed data are allowed to be rich, but in the negative theorem they
are assumed to factor through \(\pi_q\).  Examples include:
\begin{itemize}
  \item pulled-back normal state/effect values \(a\mapsto\varphi(\pi_q(a))\);
  \item represented modular data for a weight after quotienting to \(\A\);
  \item Osterwalder--Schrader or GNS data after the represented null quotient;
  \item a locally covariant functor, time-slice map, or relative Cauchy
  evolution computed only for the quotient algebra;
  \item a supplied DHR/DR sector category for \(\A\);
  \item JLMS/OAQEC relative-entropy or recovery data for a specified logical
  algebra contained in \(\A\);
  \item a continuum limit and a fixed entropy value computed only in \(\A\).
\end{itemize}

The targets are ambient.  They include full support in \(\widehat\A_q\),
faithfulness of the represented physical algebra on \(\widehat\A_q\), the
interpretation of \(\ker\pi_q\) as physical/null/gauge/inadmissible/absent,
the finite rank \(q\), and the predicate \(q=2\).  If the target is only the
represented quotient \(\A\), the negative theorem is irrelevant: standard
positive reconstruction theorems apply to the object actually supplied.

\section{Main Classification}

\begin{theorem}[support/readout and binary-rank obstruction]
Let the comparison class contain the algebras
\[
  \widehat\A_q=\A\oplus
  \bigl(\N\bar\otimes\ell^\infty(\{1,\ldots,q\})\bigr)
\]
for at least two values \(q\ge2\), including \(q=2\) and some \(q>2\).  Suppose
that every displayed datum \(D_q\) is a labelled function of the normal
quotient \(\pi_q:\widehat\A_q\to\A\), and suppose no additional datum separates
the second summand or removes it from the comparison class by a quotient,
gauge, inadmissibility, factor, split, or physical-exclusion rule.

Then all displayed data are identical as \(q\) varies.  Nevertheless
\(p_q=(0,1)\) is a nonzero central projection killed by all quotient-visible
readouts.  Consequently the data do not reconstruct:
\begin{enumerate}
  \item full ambient support in \(\widehat\A_q\);
  \item faithfulness on \(\widehat\A_q\);
  \item whether \(\ker\pi_q\) is physical, null, gauge, inadmissible, or absent;
  \item the finite central rank \(q\);
  \item the binary predicate \(q=2\).
\end{enumerate}

Moreover, if the finite central directions that survive quotienting are
encoded by \(V_q=\ell^\infty_q/\C1\), if \(N_q\subset V_q\) is the declared
null/gauge/inadmissible subspace, and if \(R_q:V_q\to\R^m\) is the available
readout map, then full finite-central accountability is exactly the rank
condition
\[
  \operatorname{rank} R_q = \dim(V_q/N_q),
\]
with \(R_q\) understood on \(V_q/N_q\).  If this accountability condition is
not met, there are non-null finite central directions that are not read out.  If
it is met, binary rank still does not follow from the Bekenstein--Hawking
coefficient alone: for every \(q\ge2\),
\[
  S_q=\log q,\qquad
  \Delta A_q=4G\hbar\log q
\]
satisfy \(S_q=\Delta A_q/(4G\hbar)\).  A binary conclusion therefore requires
an additional finite-alphabet rule selecting the least nontrivial closed
quotient, or direct \(q=2\) data with its measurement channel.
\end{theorem}

\begin{proof}
For the first part, every displayed datum has the form
\(\widetilde D\circ\pi_q\).  Hence it is independent of the second summand and
is the same for all \(q\).  The projection \(p_q\) is central and nonzero in
\(\widehat\A_q\), but \(\pi_q(p_q)=0\), so all pulled-back normal readouts and
all quotient-formed modular, sector, recovery, continuum, and entropy data
vanish on the corner \(p_q\widehat\A_qp_q\).  Replacing \(q=2\) by \(q=3\) keeps
all displayed data fixed while changing the finite central rank and the truth
value of \(q=2\).  Any rule depending only on the displayed data would assign
the same answer to both candidates, so no such rule can reconstruct the listed
ambient targets.

The accountability statement is finite linear algebra.  After quotienting by
the declared null/gauge/inadmissible directions \(N_q\), the remaining finite
central directions are \(V_q/N_q\).  A family of real-valued central readouts
separates all of them exactly when its induced map on \(V_q/N_q\) is injective.
For finite-dimensional spaces, injectivity is equivalent to
\(\operatorname{rank} R_q=\dim(V_q/N_q)\).  If the rank is smaller, a nonzero
class is both non-null and invisible to the supplied readouts.

Finally, the entropy statement is again finite.  For an accountable
\(q\)-state event alphabet, the Shannon/von Neumann entropy of the uniform
finite cell is \(\log q\).  Choosing the corresponding area increment
\(\Delta A_q=4G\hbar\log q\) gives the exact one-quarter coefficient for every
integer \(q\ge2\).  Thus the coefficient evaluates the alphabet once the
alphabet has been specified; it does not choose the alphabet.  Choosing the
least nontrivial quotient is a separate finite-alphabet rule, and a resolved
\(\log2\) spectral line is a direct datum rather than a consequence of the
quotient-visible data.
\end{proof}

\section{The Countermodel Family in Coordinates}

It is useful to spell out the same-data family without any gravitational
language.  Fix a represented algebra \(\A\), a nonzero algebra \(\N\), and a
finite rank \(q\ge2\).  The enlarged candidate is
\[
  \widehat\A_q=\A\oplus(\N\bar\otimes\ell^\infty_q).
\]
The quotient-visible experiment uses the map
\[
  \pi_q(a,n_1,\ldots,n_q)=a .
\]
Every normal state \(\omega\) on \(\A\) gives a pulled-back normal state
\(\omega\circ\pi_q\) on \(\widehat\A_q\).  Every effect, channel matrix
element, relative entropy, modular expectation value, or continuum observable
formed from the quotient is similarly pulled back.  Thus the entire displayed
experiment is independent of \(q\).

The ambient algebra, however, is not independent of \(q\).  Its center contains
the projection \(p_q=(0,1,\ldots,1)\), and the finite central quotient of the
hidden summand has \(q\) characters.  A rank-sensitive target can distinguish
\(\widehat\A_2\) from \(\widehat\A_3\), while the displayed data cannot.  This
is the whole no-go in its shortest form: same data, different target.  The
argument is insensitive to whether \(\N\) is type I, type II, or type III; if
the finite central summand is admitted and not separated or quotiented, it
survives the displayed quotient data.

There are two common ways to misread this example.  The first is to say that a
universal representation of the full algebra would see the finite center.  That
is true, but it changes the data.  The universal representation of
\(\widehat\A_q\) contains states supported on the hidden summand, while the
quotient-visible experiment does not.  The second is to say that the hidden
summand is obviously unphysical.  That may also be true in a better theory, but
then the better theory must state the physicality, quotient, gauge, or
inadmissibility rule that removes it.  The point of the countermodel is that
the quotient-visible data themselves do not provide that rule.

\subsection{Positive-support readouts versus atom separation}

For a specified finite center \(C_q=\ell^\infty_q\), a family of positive
normal readouts can be faithful on the positive cone without separating the
atoms.  In \(C_2\), the sum readout
\[
  \tau(z_1,z_2)=z_1+z_2
\]
is strictly positive on every nonzero positive vector, so it detects whether a
positive element is zero.  But it does not distinguish \((1,0)\) from
\((0,1)\), and its kernel is not an algebra ideal.  Thus positive support
faithfulness and full central tomography are different exits.  The first rules
out invisible positive support relative to the specified center; the second
reconstructs atoms, central weights, and finite rank.  A binary-rank theorem
needs the second kind of information, or a physical law replacing it.

\subsection{Static finite labels and response data}

The same finite central extension can be made static with respect to the
quotient response.  If a locally covariant or Hamiltonian response is computed
only on \(\A\), take the hidden central register to have identity evolution and
zero source response in the displayed experiment.  Then every quotient response
function is unchanged.  This does not prove that response-null finite labels
are physically admissible.  It says only that response-null labels are not
excluded by response data unless the theory also asserts response-faithful
probe admissibility, source-generated physicality, or a null/inadmissibility
quotient for response-null finite directions.

\subsection{All binary questions do not make a binary alphabet}

For \(q>2\), the finite algebra \(C_q\) has many two-outcome coarse-grained
tests:
\[
  (\chi_A,1-\chi_A),\qquad
  \varnothing\ne A\ne\{1,\ldots,q\}.
\]
Owning these yes/no questions is not the same as declaring one of them to be
the elementary formation alphabet.  A three-outcome register has three binary
questions and still has three atoms.  A four-outcome register may be a pair of
bits, a single four-letter alphabet, or a coarse quotient of a larger alphabet,
depending on which tensor/refinement structure the physical theory supplies.
Thus binary propositions are cheap; binary primitive rank is not.

\section{What Would Falsify the No-Go}

The theorem is not meant to be immune to stronger hypotheses.  It is designed
to make those hypotheses explicit.  Any of the following would block the
countermodel family.

\paragraph{Full central tomography.}
If the experiment includes normal preparations and effects separating all
minimal central projections of \(\widehat\A_q\), then \(q\) is measured.  The
same-data fiber collapses.  This is not a contradiction of the theorem because
the data no longer factor through \(\pi_q\).

\paragraph{A quotient or physicality theorem.}
If a theorem says that every finite central direction outside the represented
GNS/operational support is null, gauge, or inadmissible, then the hidden
summand is removed before ambient support is demanded.  Again, the theorem is
not contradicted; one of its excluded positive exits has been supplied.

\paragraph{Maximality of the logical algebra.}
If OAQEC/JLMS data are not merely data of a correctable subalgebra but data of
the full maximal logical algebra, then a hidden finite center in the logical
commutant is no longer invisible.  The missing input is the maximality
statement, not recovery of the supplied subalgebra.

\paragraph{A source-faithfulness principle.}
If all admitted finite labels must have nonzero source, stress, modular, or
relative-Cauchy response unless they are null, then response-null \(q\)-labels
are inadmissible.  This is a strong physical admissibility condition.  It is
plausible in some settings and false in others with transparent topological or
superselection sectors; the note treats it as an assumption that must be
proved for the class under discussion.

\paragraph{A binary finite-alphabet law.}
If every accountable finite event alphabet is required to be the least
nontrivial closed quotient of its event algebra, then the selected size is
\(2\).  The Lean core checks this finite implication.  What remains is the
physical derivation of the least-quotient rule.

\paragraph{Direct spectroscopy.}
If black-hole spectroscopy or another finite-rank probe resolves
\(\Xi=\log q\), and the measurement channel is part of the microscopic theory,
then the observed \(q\) is selected.  A result \(\Xi=\log2\) validates binary;
\(\Xi=\log3\) or \(\Xi=\log4\) falsifies it.  This route is empirical rather
than a derivation from quotient-visible data.

\section{Referee Checklist}

A referee can check the note without adopting any private formalism by asking
five questions.

\begin{enumerate}
  \item Is the comparison class explicitly closed under the finite central
  direct sums used in the countermodel?  If not, where is the physical
  exclusion theorem?
  \item Do the displayed data genuinely factor through the quotient
  \(\pi_q\)?  If some datum does not factor through \(\pi_q\), does it separate
  the hidden central rank?
  \item Is the target ambient rather than quotient-internal?  If the target is
  only the represented algebra \(\A\), the no-go is irrelevant.
  \item Is support/readout accountability supplied by full central tomography,
  by a null/gauge/inadmissibility quotient, or by the rank equation
  \(\operatorname{rank}R_q=\dim(V_q/N_q)\)?
  \item After accountability, what selects \(q=2\) rather than \(q>2\): a
  least-event law, an observed \(\log2\) spectral datum, or some other
  rank-breaking microscopic principle?
\end{enumerate}

If the answer to the fourth or fifth question is missing, the binary conclusion
is conditional.  If all five questions have sharp answers, the note provides a
map for where the positive proof should attach.

\section{How the Standard Reconstructions Enter}

\subsection{Modular and crossed-product data}

Let \(M_0\) and \(M_1\) be von Neumann algebras with faithful normal
semifinite weights \(\phi_0\) and \(\phi_1\).  For
\[
  M=M_0\oplus M_1,\qquad \phi=\phi_0\oplus\phi_1,
\]
the modular automorphism group decomposes componentwise:
\[
  \sigma_t^\phi=\sigma_t^{\phi_0}\oplus\sigma_t^{\phi_1}.
\]
Accordingly, the crossed product decomposes as
\[
  M\rtimes_{\sigma^\phi}\R
  \cong
  (M_0\rtimes_{\sigma^{\phi_0}}\R)\oplus
  (M_1\rtimes_{\sigma^{\phi_1}}\R).
\]
If the supplied modular or Type II data are formed after quotienting to
\(M_0\), they are data of \(M_0\).  They do not recover \(M_1\), its finite
central rank, or its kernel classification.  The full direct-sum core would see
the second summand, but supplying that full core is precisely an additional
non-quotient datum.  This is compatible with, and does not weaken, the
crossed-product entropy program: the entropy is an entropy of the algebra and
trace actually supplied.

\subsection{Locally covariant QFT and relative Cauchy evolution}

In the Brunetti--Fredenhagen--Verch framework a theory is a functor from
globally hyperbolic spacetimes to algebras, and relative Cauchy evolution
measures the response of that supplied functor to metric perturbations.  A
constant finite central extension
\[
  \widehat\A_q(M)=\A(M)\oplus
  \bigl(\N(M)\bar\otimes\ell^\infty_q\bigr)
\]
can be arranged so that the quotient-visible functor, its time-slice maps, and
its relative Cauchy evolution are exactly those of \(\A\).  The finite central
summand has zero quotient response because the quotient never asks it a
question.  This paragraph is not a counterexample to Fewster--Verch dynamical
locality for a fully specified theory.  It says only that the relative-Cauchy
and functorial records of the represented quotient do not decide which full
theory was admitted.  Local covariance and dynamical locality close the gap
only when imposed on the full admitted theory, or when an independent
indecomposability, source-generated closure, source faithfulness, or operational
physicality rule forbids the invisible extension.  They do not, as quotient
data alone, select \(q=2\).

\subsection{DHR/DR sector data}

DHR and Doplicher--Roberts reconstruction start from a specified observable
net and a specified category of localized transportable endomorphisms.  If the
supplied category is the category for \(\A\), and the \(q\) central characters
in an enlarged algebra are all restricted to the same object of that category,
then DR reconstruction reconstructs the field algebra and compact gauge group
associated to the supplied category.  It does not recreate central-character
fibers that were omitted from the sector datum.  The repair is also exact:
supply the amplified category \(\mathcal S(\A)\times\{1,\ldots,q\}\), or
supply another datum that separates the central characters, or prove that the
central-character fibers are null, gauge, or inadmissible.

\subsection{JLMS and operator-algebra quantum error correction}

Operator-algebra quantum error correction is intentionally algebra-relative.
Exact recovery of an algebra \(\M\) through a channel says that observables in
\(\M\) can be represented on the recovery region.  It is downward closed:
subalgebras of a correctable algebra are correctable.  Therefore exact recovery
of a diagonal or quotient algebra does not imply that the algebra is maximal
among all admitted finite central extensions.  JLMS relative-entropy equality
has the same support: it identifies the relative entropy of the supplied bulk
and boundary algebras.  To force the finite central rank, one must know that
the supplied algebra is the full logical algebra, or else supply central
tomography, central-sector calibration, or a quotient/inadmissibility rule.

\subsection{Entropy and black-hole spectroscopy}

The Bekenstein--Hawking coefficient fixes the ratio between a chosen area
increment and the entropy of the chosen finite alphabet.  It does not choose the
alphabet.  For every finite \(q\ge2\), the pair
\[
  (S_q,\Delta A_q)=(\log q,4G\hbar\log q)
\]
has the exact one-quarter coefficient.  Bekenstein--Mukhanov spectroscopy is
therefore the natural empirical boundary: a measured line spacing
\(\log q\), with the measurement channel and propagation model under control,
would select the observed \(q\).  A \(\log2\) observation would validate the
binary branch.  But that would be a new \(q\)-sensitive datum, not a derivation
from quotient-visible entropy/OS/continuum data.

\section{Positive Exits}

The theorem is useful only if it says exactly what would close the obstruction.
The current exits are:
\[
\begin{array}{p{0.34\linewidth}p{0.56\linewidth}}
\toprule
Exit & What it proves \\
\midrule
Central support faithfulness &
No nonzero positive finite central projection remains outside the represented
support. \\
\addlinespace
Full ambient central tomography &
Minimal central projections, central weights, the finite rank \(q\), and the
predicate \(q=2\) are separated by supplied readouts. \\
\addlinespace
Null/gauge/inadmissibility classifier &
Every invisible finite central direction is removed before full support is
demanded. \\
\addlinespace
Full logical/maximal correctable algebra data &
Recovery/JLMS data are known to refer to the full admitted logical algebra, not
only to a downward-closed recovered subalgebra. \\
\addlinespace
Full locally covariant/source-generated physicality &
Response-null finite labels are operationally null or inadmissible, rather
than untested physical sectors. \\
\addlinespace
Support/readout accountability rank &
For \(V_q=\ell^\infty_q/\C1\), the available readout rank equals
\(\dim(V_q/N_q)\). \\
\addlinespace
Binary finite-alphabet law &
After accountability, the elementary alphabet is the least nontrivial closed
quotient, or direct \(q=2\) data with a derived measurement channel is supplied. \\
\bottomrule
\end{array}
\]

These exits are not equivalent.  Support/readout accountability tells us that
all surviving finite central directions are either read out or removed; it does
not by itself say that the elementary event alphabet has two outcomes.  The
binary finite-alphabet law tells us how to count an accountable finite event
algebra; it does not by itself prove that all physical finite central
directions have been accounted for.  A successful positive derivation of
binary rank must supply both.

\section{Primitive Counting After Accountability}

Assume, for this section only, that the finite support problem has been solved:
the \(q\)-cell lies in the physical algebra, its central idempotents are
readable or null-quotiented, fixed-area entropy counts operationally preparable
states, and remnant-free recovery preserves the relevant input distinctions.
The binary problem still remains.

Let \(C_q=\ell^\infty(\{1,\ldots,q\})\), and suppose all binary coarse-grained
tests \((\chi_A,1-\chi_A)\) are available for nonempty proper
\(A\subset\{1,\ldots,q\}\).  This does not imply that the elementary formation
alphabet has two elements.  For \(q=3\), \(q=4\), or \(q=8\), all such binary
questions can exist while the complete \(q\)-register remains the elementary
alphabet.  Landauer cost, Blackwell/Petz sufficiency, and the area law evaluate
or compare the alphabet after the experiment has been declared; they do not
choose the experiment.

Thus a finite counting functional is needed.  Three examples are:
\[
\begin{array}{ccc}
\text{least nontrivial quotient} &\mapsto& 2,\\
\text{complete register} &\mapsto& q,\\
\text{selected } k\text{-block quotient} &\mapsto& k.
\end{array}
\]
The first rule forces binary.  The second and third do not.  The question for a
microscopic theory is therefore not whether yes/no propositions can be formed
inside a finite event algebra; they can.  The question is why the least
nontrivial quotient, rather than the complete finite register or another
coarse quotient, is the elementary counted alphabet.  In the present note this
is an assumption to be derived or measured, not a theorem.

\section{Binary Blocks Versus Elementary Binary Rank}

A recurring ambiguity is the difference between representing a finite alphabet
with binary variables and proving that the elementary physical alphabet has
two outcomes.  These are not the same statement.

If \(q=2^m\), then \(C_q\) is abstractly isomorphic to the algebra of functions
on \(m\) classical bits.  This observation is useful only after the theory has
also supplied the corresponding tensor/refinement structure as physical.  An
abstract isomorphism does not say that the \(m\) bit coordinates are available
as independently preparable, independently readable, or independently counted
subevents.  Without that extra structure, the complete \(q\)-letter register
can remain the elementary alphabet even though it admits a binary coding.

If \(q\) is not a power of two, binary source coding embeds the \(q\)-letter
alphabet into a larger binary code space with unused words.  That again does
not derive binary rank.  It replaces one physical question by another: why are
the unused code words inadmissible, null, or gauge, and why is the larger
binary code space the physical one rather than the original \(q\)-letter
alphabet?  The embedding is a representation of finite information, not a
classification of microscopic event atoms.

This distinction matters for \(q=4\) and \(q=8\).  A positive binary theory may
well reclassify them as composite binary blocks, but it can do so only when it
owns the binary factorization.  If the physical event algebra is merely a
symmetric four-letter center with no distinguished coordinate partitions, then
there is no invariant choice of two bit coordinates.  Choosing one is a
symmetry-breaking or refinement datum.  Such a datum may be physical, but it is
not supplied by the quotient-visible data.

Prime \(q>2\) is sharper.  A prime three-letter or five-letter accountable
register cannot be decomposed into smaller nontrivial Cartesian alphabets.
For those cases, a binary conclusion must either exclude the register, prove it
null/gauge/inadmissible, observe a different \(q\), or supply a rule saying
that the elementary event is a two-outcome quotient rather than the complete
prime register.  The least nontrivial closed quotient rule does exactly that,
but it is the rule to be derived.

The same issue appears in statistical sufficiency language.  Minimal
sufficient statistics are minimal relative to a specified family of future
experiments.  If the future experiment distinguishes all \(q\) outcomes, the
full \(q\)-statistic is minimal.  If the future experiment asks only a fixed
binary question, a binary statistic is sufficient.  Sufficiency theory
therefore organizes the experiment once the experiment is declared.  It does
not by itself declare which finite experiment is the microscopic formation
event.

Landauer erasure has the same limitation.  Erasing a \(k\)-letter classical
alphabet at temperature \(T\) costs \(k_B T\log k\) in the ideal reversible
bound.  The formula evaluates the alphabet.  It does not say whether the
alphabet was \(2\), \(3\), \(4\), or a binary block decomposition of \(4\).  In
the black-hole application, replacing \(\log k\) by
\(\Delta A_k/(4G\hbar)\) preserves the same structure: the one-quarter
coefficient fixes the conversion once \(k\) is known.

Thus Paper 2 should not try to prove binary by pointing to bits as a universal
language for finite information.  It should prove, assume, or observe one of
three rank-breaking inputs:
\begin{enumerate}
  \item a physical exclusion of every accountable nonbinary elementary
  register;
  \item a physical factorization/refinement theorem that demotes composite
  \(2^m\)-registers to binary blocks and excludes or quotients prime
  \(q>2\);
  \item direct \(q\)-sensitive data whose result is \(q=2\).
\end{enumerate}
Absent one of these, a binary encoding is notation, not a microscopic
derivation.

\section{Assumptions Not Proved}

\begin{longtable}{p{0.36\linewidth}p{0.54\linewidth}}
\toprule
Assumption & Status in this note \\
\midrule
The comparison class admits finite central completions
\(\widehat\A_q\) for more than one \(q\). &
This is the countermodel class.  A physical theory may forbid it, but that
requires an explicit exclusion, quotient, factor, split, or physicality rule. \\
\addlinespace
Displayed data factor through the represented quotient. &
This is the negative theorem's hypothesis.  A positive theorem must show that
the data do not factor through \(\pi_q\). \\
\addlinespace
The target is ambient support, kernel status, finite rank, or binary rank. &
If the target is only the represented quotient, ordinary positive
reconstruction theorems are untouched. \\
\addlinespace
Support/readout accountability is physically forced. &
Not proved.  The note only states the finite rank equation that would express
it. \\
\addlinespace
Operational entropy counts exactly preparable finite support. &
Not proved.  If entropy counts unpreparable labels, Page/recovery arguments do
not see them. \\
\addlinespace
The elementary finite alphabet is the least nontrivial closed quotient. &
Not proved.  This is the binary-counting law. \\
\addlinespace
Direct \(q=2\) spectroscopy exists and its measurement channel is derived. &
Not supplied.  Such data would validate binary empirically but would not be a
derivation from quotient-visible data. \\
\bottomrule
\end{longtable}

\section{Lean Core}

The accompanying Lean archive
\[
  \texttt{docs/support\_readout\_q2\_lean\_verification/}
\]
checks only the finite logical core.  It imports Lean's \texttt{Std} library
and proves:
\begin{itemize}
  \item a same-data pair with different target values forbids reconstruction
  from the data;
  \item the same statement for a proposition-valued target;
  \item the blind \(q=2\) and \(q=3\) candidates have the same data but
  different rank and different truth values of ``rank is two'';
  \item a missing centered readout/null rank prevents accountability;
  \item if the least nontrivial event-size rule is supplied for \(q\ge2\), it
  selects size \(2\).
\end{itemize}
It does not formalize von Neumann algebras, DHR theory, locally covariant QFT,
JLMS, crossed products, or black-hole spectroscopy.  Those are literature
inputs to the manuscript.  The Lean file is therefore evidence for the finite
descent kernel, not evidence for the analytic or physical assumptions.

\section{Conclusion}

The support/readout obstruction is not a failure of AQFT or quantum
information reconstruction.  It is a boundary on what the displayed data are
allowed to imply.  If all data factor through a represented quotient, they
cannot decide the status or rank of an invisible finite central kernel.  If the
finite center is made accountable, entropy still does not choose \(q=2\);
binary rank needs a finite-alphabet law or direct binary data.

This turns the remaining binary problem into a short checklist.  A positive
theory must either derive support/readout accountability and then derive the
least-nontrivial-alphabet law, or it must produce a \(q\)-sensitive measurement
whose channel is itself part of the theory and whose outcome is \(\log2\).
Without one of those routes, the honest conclusion is conditional framework,
not primitive unavoidability.

\begin{thebibliography}{99}

\bibitem{BFV}
R.~Brunetti, K.~Fredenhagen, and R.~Verch,
``The generally covariant locality principle: A new paradigm for local quantum
physics,''
Commun. Math. Phys. \textbf{237} (2003), 31--68,
arXiv:math-ph/0112041.

\bibitem{FewsterVerch}
C.~J. Fewster and R.~Verch,
``Dynamical locality and covariance: What makes a physical theory the same in
all spacetimes?''
Ann. Henri Poincare \textbf{13} (2012), 1613--1674,
arXiv:1106.4785.

\bibitem{DHR1}
S.~Doplicher, R.~Haag, and J.~E. Roberts,
``Local observables and particle statistics I,''
Commun. Math. Phys. \textbf{23} (1971), 199--230.

\bibitem{DHR2}
S.~Doplicher, R.~Haag, and J.~E. Roberts,
``Local observables and particle statistics II,''
Commun. Math. Phys. \textbf{35} (1974), 49--85.

\bibitem{DR}
S.~Doplicher and J.~E. Roberts,
``Why there is a field algebra with a compact gauge group describing the
superselection structure in particle physics,''
Commun. Math. Phys. \textbf{131} (1990), 51--107.

\bibitem{JLMS}
D.~L. Jafferis, A.~Lewkowycz, J.~Maldacena, and S.~J. Suh,
``Relative entropy equals bulk relative entropy,''
JHEP \textbf{06} (2016), 004,
arXiv:1512.06431.

\bibitem{HarlowRT}
D.~Harlow,
``The Ryu--Takayanagi Formula from Quantum Error Correction,''
Commun. Math. Phys. \textbf{354} (2017), 865--912,
arXiv:1607.03901.

\bibitem{DHW}
X.~Dong, D.~Harlow, and A.~C. Wall,
``Reconstruction of Bulk Operators within the Entanglement Wedge in
Gauge-Gravity Duality,''
Phys. Rev. Lett. \textbf{117} (2016), 021601,
arXiv:1601.05416.

\bibitem{KangKolchmeyer}
M.~J. Kang and D.~K. Kolchmeyer,
``Holographic Relative Entropy in Infinite-dimensional Hilbert Spaces,''
Commun. Math. Phys. (2023),
arXiv:1811.05482.

\bibitem{WittenCrossed}
E.~Witten,
``Gravity and the Crossed Product,''
JHEP \textbf{10} (2022), 008,
arXiv:2112.12828.

\bibitem{CLPW}
V.~Chandrasekaran, R.~Longo, G.~Penington, and E.~Witten,
``An Algebra of Observables for de Sitter Space,''
JHEP \textbf{02} (2023), 082,
arXiv:2206.10780.

\bibitem{CPW}
V.~Chandrasekaran, G.~Penington, and E.~Witten,
``Large \(N\) Algebras and Generalized Entropy,''
JHEP \textbf{04} (2023), 009,
arXiv:2209.10454.

\bibitem{BekensteinMukhanov}
J.~D. Bekenstein and V.~F. Mukhanov,
``Spectroscopy of the quantum black hole,''
Phys. Lett. B \textbf{360} (1995), 7--12,
arXiv:gr-qc/9505012.

\end{thebibliography}

\end{document}
